\documentclass[10pt,a4paper]{amsart}
\usepackage{amsmath,amssymb,amsthm,graphicx}
% Retained preamble for the Stacks Project chapter snapshot in this directory.
% The chapter's own preamble.tex is a copy of the shared Stacks Project preamble; the block
% below is sliced MECHANICALLY, line for line, from it. Nothing was added, removed or
% redefined; the wrapper class, the xr-hyper/hyperref wiring and the externaldocument list
% are omitted because the assembler supplies the document class.
%
% lines taken from preamble.tex: 3, 5, 9, 11, 12, 156, 157, 158, ... (51 lines total)
%
% package coverage restored from the chapter's own preamble (2026-10-13): xy, xy-2cell

\usepackage[all]{xy}
\xyoption{2cell}
\UseAllTwocells

\newenvironment{reference}{\comment}{\endcomment}
\newenvironment{slogan}{\comment}{\endcomment}
\newenvironment{history}{\comment}{\endcomment}
\theoremstyle{plain}
\newtheorem{theorem}[subsection]{Theorem}
\newtheorem{proposition}[subsection]{Proposition}
\newtheorem{lemma}[subsection]{Lemma}
\theoremstyle{definition}
\newtheorem{definition}[subsection]{Definition}
\newtheorem{example}[subsection]{Example}
\newtheorem{exercise}[subsection]{Exercise}
\newtheorem{situation}[subsection]{Situation}
\theoremstyle{remark}
\newtheorem{remark}[subsection]{Remark}
\newtheorem{remarks}[subsection]{Remarks}
\numberwithin{equation}{subsection}
\def\lim{\mathop{\mathrm{lim}}\nolimits}
\def\colim{\mathop{\mathrm{colim}}\nolimits}
\def\Spec{\mathop{\mathrm{Spec}}}
\def\Hom{\mathop{\mathrm{Hom}}\nolimits}
\def\Ext{\mathop{\mathrm{Ext}}\nolimits}
\def\SheafHom{\mathop{\mathcal{H}\!\mathit{om}}\nolimits}
\def\SheafExt{\mathop{\mathcal{E}\!\mathit{xt}}\nolimits}
\def\Sch{\mathit{Sch}}
\def\Mor{\mathop{\mathrm{Mor}}\nolimits}
\def\Ob{\mathop{\mathrm{Ob}}\nolimits}
\def\Sh{\mathop{\mathit{Sh}}\nolimits}
\def\NL{\mathop{N\!L}\nolimits}
\def\CH{\mathop{\mathrm{CH}}\nolimits}
\def\proetale{{pro\text{-}\acute{e}tale}}
\def\etale{{\acute{e}tale}}
\def\QCoh{\mathit{QCoh}}
\def\Ker{\mathop{\mathrm{Ker}}}
\def\Im{\mathop{\mathrm{Im}}}
\def\Coker{\mathop{\mathrm{Coker}}}
\def\Coim{\mathop{\mathrm{Coim}}}
\DeclareMathSymbol{\boxtimes}{\mathbin}{AMSa}{"02}
\def\QCohstack{\mathcal{QC}\!\mathit{oh}}
\def\Cohstack{\mathcal{C}\!\mathit{oh}}
\def\Spacesstack{\mathcal{S}\!\mathit{paces}}
\def\Quotfunctor{\mathrm{Quot}}
\def\Hilbfunctor{\mathrm{Hilb}}
\def\Curvesstack{\mathcal{C}\!\mathit{urves}}
\def\Polarizedstack{\mathcal{P}\!\mathit{olarized}}
\def\Complexesstack{\mathcal{C}\!\mathit{omplexes}}
\def\Pic{\mathop{\mathrm{Pic}}\nolimits}
\def\Picardstack{\mathcal{P}\!\mathit{ic}}
\def\Picardfunctor{\mathrm{Pic}}
\def\Deformationcategory{\mathcal{D}\!\mathit{ef}}
\usepackage{cleveref}
\crefname{originallemma}{original lemma}{original lemmas}
\Crefname{originallemma}{Original lemma}{Original lemmas}
\makeatletter
\newlabel{categories-lemma-yoneda}{{001P}{source}{}{sourcepointer-1}{}}
\newlabel{categories-lemma-yoneda@cref}{{[originallemma][0][]001P}{[1][0][]source}{}{}{}}
\makeatother
\crefname{originallemma}{original lemma}{original lemmas}
\Crefname{originallemma}{Original lemma}{Original lemmas}
\makeatletter
\newlabel{divisors-lemma-fibre-Cartier}{{062Y}{source}{}{sourcepointer-2}{}}
\newlabel{divisors-lemma-fibre-Cartier@cref}{{[originallemma][0][]062Y}{[1][0][]source}{}{}{}}
\makeatother
\crefname{originallemma}{original lemma}{original lemmas}
\Crefname{originallemma}{Original lemma}{Original lemmas}
\makeatletter
\newlabel{lemma-define-open}{{0B9V}{source}{}{sourcepointer-3}{}}
\newlabel{lemma-define-open@cref}{{[originallemma][0][]0B9V}{[1][0][]source}{}{}{}}
\makeatother
\crefname{originallemma}{original lemma}{original lemmas}
\Crefname{originallemma}{Original lemma}{Original lemmas}
\makeatletter
\newlabel{lemma-divisors-on-curves}{{0B9D}{source}{}{sourcepointer-4}{}}
\newlabel{lemma-divisors-on-curves@cref}{{[originallemma][0][]0B9D}{[1][0][]source}{}{}{}}
\makeatother
\crefname{originallemma}{original lemma}{original lemmas}
\Crefname{originallemma}{Original lemma}{Original lemmas}
\makeatletter
\newlabel{perfect-lemma-flat-proper-perfect-direct-image-general}{{0B91}{source}{}{sourcepointer-5}{}}
\newlabel{perfect-lemma-flat-proper-perfect-direct-image-general@cref}{{[originallemma][0][]0B91}{[1][0][]source}{}{}{}}
\makeatother
\crefname{originallemma}{original lemma}{original lemmas}
\Crefname{originallemma}{Original lemma}{Original lemmas}
\makeatletter
\newlabel{perfect-lemma-open-where-cohomology-in-degree-i-rank-r-geometric}{{0B9S}{source}{}{sourcepointer-6}{}}
\newlabel{perfect-lemma-open-where-cohomology-in-degree-i-rank-r-geometric@cref}{{[originallemma][0][]0B9S}{[1][0][]source}{}{}{}}
\makeatother
\crefname{originallemma}{original lemma}{original lemmas}
\Crefname{originallemma}{Original lemma}{Original lemmas}
\makeatletter
\newlabel{proposition-hilb-d}{{0B9I}{source}{}{sourcepointer-7}{}}
\newlabel{proposition-hilb-d@cref}{{[originallemma][0][]0B9I}{[1][0][]source}{}{}{}}
\makeatother
\crefname{originallemma}{original lemma}{original lemmas}
\Crefname{originallemma}{Original lemma}{Original lemmas}
\makeatletter
\newlabel{remark-universal-object-hilb-d}{{0B9J}{source}{}{sourcepointer-8}{}}
\newlabel{remark-universal-object-hilb-d@cref}{{[originallemma][0][]0B9J}{[1][0][]source}{}{}{}}
\makeatother
\crefname{originallemma}{original lemma}{original lemmas}
\Crefname{originallemma}{Original lemma}{Original lemmas}
\makeatletter
\newlabel{varieties-definition-degree-invertible-sheaf}{{0AYR}{source}{}{sourcepointer-9}{}}
\newlabel{varieties-definition-degree-invertible-sheaf@cref}{{[originallemma][0][]0AYR}{[1][0][]source}{}{}{}}
\makeatother
\title{The open subfunctor F is representable by an open of Hilb\textasciicircum{}g}
\author{The Stacks Project Authors}
\date{stacks-project:pic.tex:local-snapshot-20261012; GNU FDL 1.2 or later}
\begin{document}
\maketitle
\noindent\textit{Editorial note.} Statement span is lines 1213-1219 of sources/stacks-pic/pinned-original-source.tex; proof span is lines 1221-1314, the complete author proof (92 body lines).\par\medskip
\noindent\textit{Editorial note.} Navigation stubs: 9 cited-but-not-included labels are supplied as non-load-bearing pointers with their pinned tags where the snapshot tag index has them.\par\medskip
\noindent\textit{Editorial note.} extra\_packages is empty on purpose: the retained preamble already loads xy with options and a second declaration clashes in the assembler wrapper.\par\medskip
\noindent\textit{Editorial note.} Text is retained verbatim, including the author's own typos and placeholders.\par\medskip
\noindent\textit{Editorial note.} Queue index 2 of queue\_pic.json (pb=3806, stmt 1213-1219, proof 1221-1314).\par\medskip
\noindent\textit{Editorial note.} stacks-pic wave-1 (pi01): fresh adjudication accept (2026-10-14); pool per standing queue hypercovering -> pic (SMC reassigned to Zhen per Chief 08:5x); pic.tex 27 labels/23 with proof judged 21 accept / 2 reject / 4 pb=0 env-dup rows excluded; spans cite-free and reference/slogan/history-free so no patchenv needed. Named omission: Author omits the verification that the assignment L |-> m is an inverse transformation, and leaves 'a nonzero section defines an effective Cartier divisor' as clear - named mechanical checks; the W construction and the divisor claim are in-span.\par\medskip
\noindent\textit{Allowed standard prerequisites (editorial).}\begin{itemize}
\item Explicit prerequisite boundary. The cited inputs are prerequisites; the argument in the retained span is not replaced by any cited result.
\end{itemize}\medskip
\ifdefined\hypertarget\hypertarget{sourcepointer-1}{}\fi
\noindent\textit{Source reference.} Original lemma 001P: sources/stacks-pic/pinned-original-source.tex, source lines 1314--1314. Cited by the retained statement or proof but not part of the collected spans; supplied only as a navigation pointer. The argument does not depend on its content.\par\medskip
\ifdefined\hypertarget\hypertarget{sourcepointer-2}{}\fi
\noindent\textit{Source reference.} Original lemma 062Y: sources/stacks-pic/pinned-original-source.tex, source lines 1314--1314. Cited by the retained statement or proof but not part of the collected spans; supplied only as a navigation pointer. The argument does not depend on its content.\par\medskip
\ifdefined\hypertarget\hypertarget{sourcepointer-3}{}\fi
\noindent\textit{Source reference.} Original lemma 0B9V: sources/stacks-pic/pinned-original-source.tex, source lines 1187--1195. Cited by the retained statement or proof but not part of the collected spans; supplied only as a navigation pointer. The argument does not depend on its content.\par\medskip
\ifdefined\hypertarget\hypertarget{sourcepointer-4}{}\fi
\noindent\textit{Source reference.} Original lemma 0B9D: sources/stacks-pic/pinned-original-source.tex, source lines 625--634. Cited by the retained statement or proof but not part of the collected spans; supplied only as a navigation pointer. The argument does not depend on its content.\par\medskip
\ifdefined\hypertarget\hypertarget{sourcepointer-5}{}\fi
\noindent\textit{Source reference.} Original lemma 0B91: sources/stacks-pic/pinned-original-source.tex, source lines 1314--1314. Cited by the retained statement or proof but not part of the collected spans; supplied only as a navigation pointer. The argument does not depend on its content.\par\medskip
\ifdefined\hypertarget\hypertarget{sourcepointer-6}{}\fi
\noindent\textit{Source reference.} Original lemma 0B9S: sources/stacks-pic/pinned-original-source.tex, source lines 1314--1314. Cited by the retained statement or proof but not part of the collected spans; supplied only as a navigation pointer. The argument does not depend on its content.\par\medskip
\ifdefined\hypertarget\hypertarget{sourcepointer-7}{}\fi
\noindent\textit{Source reference.} Original lemma 0B9I: sources/stacks-pic/pinned-original-source.tex, source lines 849--867. Cited by the retained statement or proof but not part of the collected spans; supplied only as a navigation pointer. The argument does not depend on its content.\par\medskip
\ifdefined\hypertarget\hypertarget{sourcepointer-8}{}\fi
\noindent\textit{Source reference.} Original lemma 0B9J: sources/stacks-pic/pinned-original-source.tex, source lines 912--931. Cited by the retained statement or proof but not part of the collected spans; supplied only as a navigation pointer. The argument does not depend on its content.\par\medskip
\ifdefined\hypertarget\hypertarget{sourcepointer-9}{}\fi
\noindent\textit{Source reference.} Original lemma 0AYR: sources/stacks-pic/pinned-original-source.tex, source lines 1314--1314. Cited by the retained statement or proof but not part of the collected spans; supplied only as a navigation pointer. The argument does not depend on its content.\par\medskip
\begin{lemma}
\label{lemma-open-representable}
Let $k$ be a field. Let $X$ be a smooth projective curve of genus $g$
over $k$ with a $k$-rational point $\sigma$. The open subfunctor $F$ defined
in Lemma \ref{lemma-define-open} is representable by an open subscheme of
$\underline{\Hilbfunctor}^g_{X/k}$.
\end{lemma}

\begin{proof}
In this proof unadorned products are over $\Spec(k)$.
By Proposition \ref{proposition-hilb-d} the scheme
$H = \underline{\Hilbfunctor}^g_{X/k}$ exists.
Consider the universal divisor $D_{univ} \subset H \times X$
and the associated invertible sheaf $\mathcal{O}(D_{univ})$, see
Remark \ref{remark-universal-object-hilb-d}.
We adjust by tensoring with the pullback via
$\sigma_H : H \to H \times X$ to get
$$
\mathcal{L}_H =
\mathcal{O}(D_{univ})
\otimes_{\mathcal{O}_{H \times X}}
\text{pr}_H^*\sigma_H^*\mathcal{O}(D_{univ})^{\otimes -1}
\in
\Picardfunctor_{X/k, \sigma}(H)
$$
By the Yoneda lemma (Categories, Lemma \ref{categories-lemma-yoneda})
the invertible sheaf $\mathcal{L}_H$ defines a natural transformation
$$
h_H \longrightarrow \Picardfunctor_{X/k, \sigma}
$$
Because $F$ is an open subfuctor, there exists a maximal open
$W \subset H$ such that $\mathcal{L}_H|_{W \times X}$ is in
$F(W)$. Of course, this open is nothing else than the
open subscheme constructed in
Derived Categories of Schemes, Lemma
\ref{perfect-lemma-open-where-cohomology-in-degree-i-rank-r-geometric}
with $i = 0$ and $r = 1$ for the morphism $H \times X \to H$ and the sheaf
$\mathcal{F} = \mathcal{O}(D_{univ})$. Applying the Yoneda lemma
again we obtain a commutative diagram
$$
\xymatrix{
h_W \ar[d] \ar[r] & F \ar[d] \\
h_H \ar[r] & \Picardfunctor_{X/k, \sigma}
}
$$
To finish the proof we will show that the top horizontal arrow is an
isomorphism.

\medskip\noindent
Let $\mathcal{L} \in F(T) \subset \Picardfunctor_{X/k, \sigma}(T)$.
Let $\mathcal{N}$ be the invertible $\mathcal{O}_T$-module
such that $Rf_{T, *}\mathcal{L} \cong \mathcal{N}[0]$.
The adjunction map
$$
f_T^*\mathcal{N} \longrightarrow \mathcal{L}
\quad\text{corresponds to a section }s\text{ of}\quad
\mathcal{L} \otimes f_T^*\mathcal{N}^{\otimes -1}
$$
on $X_T$. Claim: The zero scheme of $s$ is a relative effective Cartier
divisor $D$ on $(T \times X)/T$ finite locally free of degree $g$ over $T$.

\medskip\noindent
Let us finish the proof of the lemma admitting the claim.
Namely, $D$ defines a morphism $m : T \to H$ such that $D$ is the pullback of
$D_{univ}$. Then
$$
(m \times \text{id}_X)^*\mathcal{O}(D_{univ}) \cong
\mathcal{O}_{T \times X}(D)
$$
Since $\mathcal{O}_{T \times X}(D) \cong
\mathcal{L} \otimes f_T^*\mathcal{N}^{\otimes -1}$
we conclude that $(m \times \text{id}_X)^*\mathcal{L}_H$ and $\mathcal{L}$
differ by the pullback of an invertible module on $T$.
This shows that $m : T \to H$ factors through the open $W \subset H$ above.
Moreover, it follows that these invertible modules define, after adjusting
by pullback via $\sigma_T$ as above, the same element of
$\Picardfunctor_{X/k, \sigma}(T)$. Chasing diagrams using Yoneda's lemma
we see that $m \in h_W(T)$ maps to $\mathcal{L} \in F(T)$. We omit
the verification that the rule $F(T) \to h_W(T)$,
$\mathcal{L} \mapsto m$ defines an inverse of the transformation
of functors above.

\medskip\noindent
Proof of the claim. Since $D$ is a locally principal closed subscheme
of $T \times X$, it suffices to show that the fibres of $D$ over $T$ are
effective Cartier divisors, see Lemma \ref{lemma-divisors-on-curves} and
Divisors, Lemma \ref{divisors-lemma-fibre-Cartier}. Because taking
cohomology of $\mathcal{L}$ commutes with base change
(Derived Categories of Schemes, Lemma
\ref{perfect-lemma-flat-proper-perfect-direct-image-general})
we reduce to $T = \Spec(K)$ where $K/k$ is a field extension.
Then $\mathcal{L}$ is an invertible sheaf on $X_K$ with
$H^0(X_K, \mathcal{L}) = K$ and $H^1(X_K, \mathcal{L}) = 0$. Thus
$$
\deg(\mathcal{L}) = \chi(X_K, \mathcal{L}) - \chi(X_K, \mathcal{O}_{X_K})
= 1 - (1 - g) = g
$$
See Varieties, Definition \ref{varieties-definition-degree-invertible-sheaf}.
To finish the proof we have to show a nonzero section of $\mathcal{L}$
defines an effective Cartier divisor on $X_K$.
This is clear.
\end{proof}

\end{document}
